\section{Introduction}

WiFi Channel State Information (CSI) sensing---activity recognition, gesture detection, localization using wireless signals---is a growing field with a well-documented deployment problem: models trained in one room fail in another, with accuracy drops of 30--50\% commonly reported~\cite{yousefi2017uthar}. Test-time adaptation (TTA) is the natural solution: WiFi infrastructure provides continuous unlabeled data from the target environment, exactly matching TTA's problem formulation.

\textbf{But does TTA actually work for WiFi CSI?} This question has not been systematically studied. Prior WiFi sensing work focuses on supervised domain adaptation or domain generalization, skipping the label-free deployment-time regime entirely. Meanwhile, TTA research evaluates almost exclusively on vision corruption benchmarks (ImageNet-C, CIFAR-C), where shifts are statistical rather than physical.

We fill this gap with \textbf{WiFi-TTA-Bench}, a WiFi CSI benchmark designed to answer whether TTA methods that succeed on vision benchmarks also succeed on \emph{physics-structured} domain shifts---shifts caused by changes in the physical measurement process (multipath propagation) rather than surface-level corruptions.

\paragraph{Headline counterexample.} On Widar's hard cross-room shift, faithful BN-only TENT on an MLP+BN backbone averages $-$0.61~pp ($n{=}15$, Table~\ref{tab:faithful_tent}). On SignFi's milder temporal split, the \emph{same} BN-only TENT on the \emph{same} MLP+BN backbone averages $+$21.1~pp ($n{=}15$, $95\%\text{CI}=[+13.1,+29.6]$, Table~\ref{tab:signfi_mlpbn}). On Widar, no method evaluated on the primary MLP backbone has a 95\% CI clearing zero, and our simplified shared-MLP ports of SAR/CoTTA show strongly negative mean gain on all three real datasets; their faithful BN-only counterparts shrink toward zero. This pattern suggests that, on the evaluated WiFi CSI splits, the locus of TTA outcome depends on the interaction between shift severity, architecture, and source calibration rather than on the choice of adaptation method alone.

\paragraph{Evaluative role.} Following the NeurIPS 2026 Evaluations \& Datasets track framing, WiFi-TTA-Bench is first an evaluation artifact rather than a method contribution. We state its evaluative role explicitly: the benchmark tests the hypothesis that existing TTA methods---evaluated to date almost exclusively on vision corruption benchmarks---transfer to WiFi CSI under environmental and temporal shifts. It supports claims of the form ``under leave-one-environment-out splits on \{Widar\_BVP, NTU-Fi, SignFi-10\} with the shared MLP or MLP+BN backbones, method $M$ achieves mean gain $g$ with CI $[l,u]$ and negative-adaptation rate $r$.'' Its scope does not extend to cross-modality generalisation, device-in-the-loop latency, or real-time deployment. Its primary limitation is that it does not collect new raw CSI and therefore inherits the three upstream datasets' scale and split characteristics (two of three are temporal splits within a single laboratory).

\paragraph{Contributions.}
\begin{enumerate}
    \item \textbf{Empirical finding: evidence suggests a shift-severity $\times$ architecture interaction on WiFi CSI.} The SignFi/Widar counterexample above rules out a universal ``TTA fails on physics-structured shift'' interpretation while still supporting the weaker claim that no evaluated method achieves a positive CI on the primary MLP backbone on any real dataset.
    \item \textbf{Harm-aware benchmark} with 12 methods on Widar (10--12 on NTU-Fi/SignFi; coverage matrix in Appendix~\ref{app:coverage}), 3 real datasets, synthetic control, 15 paired observations per method, explicit fidelity labels (TENT-proj vs faithful BN-only TENT; SHOT as IM; simplified SAR/CoTTA/LAME), and a 5-metric protocol (gain, CI, Cohen's $d$, negative-adaptation rate, source drop).
    \item \textbf{Diagnostic evidence consistent with three failure modes}: (i) overconfident target predictions---source-model mean max-softmax confidence is 70--81\,\% across every Widar class while per-class accuracies range from 24\,\% to 57\,\% (Appendix Table~\ref{tab:overconfidence}); (ii) batch-statistic unreliability at WiFi-scale batches (faithful BN-only TENT $-$0.61~pp, BN-reset $-$1.11~pp on Widar); (iii) a large synthetic-to-real swing for the physics-prior objective family ($+$7.4 to $+$8.0 on synthetic $\to$ $-$0.5~pp on real Widar). Label smoothing ($\alpha{=}0.1$) on Widar gives a small, non-significant harm reduction (negative-adaptation rate 0.47 $\to$ 0.27 for the best gradient-based method; Appendix~\ref{app:source_overfit}); a labelled-target oracle upper bound ($+$24.9~pp mean gap over 3 rooms $\times$ 3 seeds, \texttt{scripts/train\_supervised\_upper\_bound.py}; Appendix~\ref{app:upperbound}) quantifies that unlabelled TTA under this protocol recovers $\approx$1\% of the available headroom.
    \item \textbf{Reproducible artifact}: code at \url{https://anonymous.4open.science/r/wifi-tta-bench-ed2026}, mirrored Croissant 1.0 dataset card at \url{https://huggingface.co/datasets/anonymous-ed2026/wifi-tta-bench}, data-preparation scripts, split manifests, a fail-fast reproducer (\texttt{scripts/reproduce\_all\_results.py}), and 248 automated tests (both URLs anonymised for double-blind review).
\end{enumerate}
